\documentclass[a4paper,12pt]{article}
\usepackage{graphicx}
\usepackage[a4paper, top=0cm, bottom=2cm, left=2cm, right=2cm]{geometry} % Ajuste les marges
\usepackage{xcolor} % Pour ajouter des couleurs
\usepackage{hyperref} % Pour avoir des références colorées si nécessaire
\usepackage{eso-pic}         % Pour ajouter des éléments en arrière-plan

\usepackage[french]{babel}
\usepackage[T1]{fontenc}
\usepackage{mathrsfs}
\usepackage{amsmath}
\usepackage{amsfonts}
\usepackage{amssymb}
\usepackage{tkz-tab}

\usepackage{tikz}
\usetikzlibrary{arrows, shapes.geometric, fit}
\newcounter{correction} % Compteur pour les questions

% Définir la commande pour afficher une question numérotée
\newcommand{\question}{%
  \refstepcounter{correction}%
  \textbf{\textcolor{black}{Question \thecorrection :}} \ignorespaces
}
% Commande pour ajouter du texte en arrière-plan
\AddToShipoutPicture{
    \AtTextCenter{%
        \makebox[0pt]{\rotatebox{80}{\textcolor[gray]{0.6}{\fontsize{10cm}{10cm}\selectfont PGB}}}
    }
}
\begin{document}
\hrule % Barre horizontale
% En-tête
\begin{center}
    \begin{tabular}{@{} p{5cm} p{5cm} p{5cm} @{}} % 3 colonnes avec largeurs fixées
        Lycée Dindéfélo & \quad\quad\quad Test 4 & 29 Novembre 2025 \\
    \end{tabular}
    \\[-0.01cm] % Ajuster l'espace vertical entre le tableau et la barre
    \hrule % Barre horizontale
\end{center}
\begin{center}
    \textbf{\Large Calcul Littéral } \\[0.2cm]
    \textbf{\large Professeur : M. BA} \\[0.2cm]
    \textbf{Classe : 4 M2} \\[0.2cm]
    \textbf{\small Durée : 10 minutes} \\[0.2cm]
    \textbf{\small Note :\quad\quad /5}
\end{center}

% Nom de l'élève
\textbf{\small Nom de l'élève :} \underline{\hspace{8cm}} \\[0.5cm]

\question (1,5 points)\\
\[
\begin{aligned}
A&=-2-6+2+6\\
&=\\
&=\\
&=\\
\end{aligned}
\quad\quad\quad
\begin{aligned}
B&=-12+11-8+12\\
&=\\
&=\\
&=\\
\end{aligned}
\quad\quad\quad
\begin{aligned}
C&=-9-15+9-21\\
&=\\
&=\\
&=\\
\end{aligned}
\]
\[
\begin{aligned}
D&= 12 + 3 \times 5\\
&=\\
&=\\
&=\\
&=\\
\end{aligned}
\quad\quad\quad
\begin{aligned}
E&= 40 - (8 + 3 \times 2)\\
&=\\
&=\\
&=\\
&=\\
\end{aligned}
\quad\quad\quad
\begin{aligned}
F&= 40 - [8 + 3 \times (2 -  12 + 3) \times 5]\\
&=\\
&=\\
&=\\
&=\\
\end{aligned}
\]
\\[1cm]
\question (1,5 points)\\
\[
\begin{aligned}
D&=(-2)\times(-7)+5\\
&=\\
&=\\
&\\
&\\
\end{aligned}
\quad\quad\quad
\begin{aligned}
E&=-15-(5)\times(-9)+9\\
&=\\
&=\\
&=\\
&\\
\end{aligned}
\quad\quad
\begin{aligned}
F&=2-(-11)\times(7)-15+2-(-10)\times(4)\\
&=\\
&=\\
&=\\
&=\\
\end{aligned}
\]
\\[1cm]
\question (2 points)\\
\[\begin{aligned}
G &= -\frac{7}{5} + \frac{9}{10} - \frac{12}{5} + \frac{4}{10} \\
   &= \\
   &= \\
   &= \\
   &= \\
\end{aligned} 
\quad\quad\quad 
\begin{aligned}
H &= -\frac{7}{3} \times \frac{5}{2} + \frac{35}{4} \times \frac{4}{6} \\
   &= \\
   &= \\
   &= \\
    &= \\
\end{aligned}
\]
\\[1cm]
\question (2 points)\\
+++++++++++++++++++++++++++

\section*{Exercice 2 : Parenthèses et crochets (Niveau 5ème)}
\textit{Calculer les expressions suivantes en respectant les priorités (parenthèses, puis crochets).}

\begin{enumerate}
    \item $E = 15 - [7 - (3 + 1)]$
    \item $F = 3 \times [12 - (5 - 2)]$
    \item $G = \dfrac{18 + 6}{5 - 2}$
\end{enumerate}

\vspace{1cm}

\section*{Exercice 3 : Nombres relatifs et règles de signes (Niveau 4ème)}
\textit{Effectuer les calculs suivants en appliquant les règles de signes.}

\begin{enumerate}
    \item $H = (-8) + (+5) - (-3) + (-12)$
    \item $I = (-4) \times (+6)$
    \item $J = (-5) \times (-3) \times (-2)$
    \item $K = \dfrac{-24}{-4}$
    \item $L = (+15) - [(-7) - (+3) \times (-2)]$
\end{enumerate}

\vspace{1cm}

\section*{Exercice 4 : Bilan et complexité (Niveau 4ème / Approfondissement)}
\textit{Calculer l'expression suivante en détaillant chaque étape :}

$$M = -10 + \frac{15 - 3 \times 5}{-2} - [(-4) \times (-2)]$$
\end{document}  